\documentclass[preview]{standalone}

\usepackage{amsmath}
\usepackage{amssymb}
\usepackage{stellar}
\usepackage{definitions}
\usepackage{bettelini}

\begin{document}

\title{Stellar}
\id{linearalgebra-orientation}
\genpage

\section{Euclidean orientation}

\begin{snippet}{determinant-area-orientation-r2}
    Let \(v_1,v_2\in\realnumbers^2\), and let
    \[
        A=
        \begin{pmatrix}
            \vert & \vert\\
            v_1 & v_2\\
            \vert & \vert
        \end{pmatrix}.
    \]
    Then \(|\det(A)|\) is the area of the parallelogram spanned by
    \(v_1\) and \(v_2\). Moreover, the sign of \(\det(A)\) distinguishes the
    two possible orientations of the plane: positive \determinant means that
    the ordered pair \((v_1,v_2)\) has the same \orientation as the canonical
    ordered \basis, while negative \determinant means that it has the opposite
    \orientation.
\end{snippet}

\begin{snippetexample}{rotation-preserves-orientation-example}{Rotations preserve orientation}
    The \matrix of the rotation of angle \(\theta\) in the canonical
    \basis of \(\realnumbers^2\) is
    \[
        R_\theta=
        \begin{pmatrix}
            \cos\theta & -\sin\theta\\
            \sin\theta & \cos\theta
        \end{pmatrix}.
    \]
    Its determinant is
    \[
        \det(R_\theta)
        =
        \cos^2\theta+\sin^2\theta
        =
        1.
    \]
    Therefore rotations preserve \orientation.
\end{snippetexample}

\begin{snippetexample}{reflection-reverses-orientation-example}{A reflection reverses orientation}
    The reflection of \(\realnumbers^2\) across the vertical axis sends
    \[
        (1,0)\mapsto(-1,0),
        \qquad
        (0,1)\mapsto(0,1).
    \]
    Its \matrix in the canonical \basis is
    \[
        S=
        \begin{pmatrix}
            -1 & 0\\
            0 & 1
        \end{pmatrix}.
    \]
    Since
    \[
        \det(S)=-1,
    \]
    this reflection reverses \orientation.
\end{snippetexample}

\section{Automorphisms}

\begin{snippetdefinition}{linear-automorphism-definition}{Linear automorphism}
    Let \(V\) be a \vectorspace over a \field \(\mathbb{F}\). A
    \emph{linear automorphism} of \(V\) is an invertible
    \lineartransformation
    \[
        f\colon V\fromto V.
    \]
    The \set of all \linearautomorphism[linear automorphisms] of \(V\) is denoted
    \[
        \operatorname{Aut}(V).
    \]
\end{snippetdefinition}

\begin{snippetdefinition}{orientation-preserving-automorphism-definition}{Orientation-preserving automorphism}
    Let \(V\) be a finite-dimensional real \vectorspace, and let
    \(f\in\operatorname{Aut}(V)\). The \linearautomorphism \(f\) is
    \emph{orientation-preserving} if
    \[
        \det(f)>0.
    \]
\end{snippetdefinition}

\begin{snippetdefinition}{orientation-reversing-automorphism-definition}{Orientation-reversing automorphism}
    Let \(V\) be a finite-dimensional real \vectorspace, and let
    \(f\in\operatorname{Aut}(V)\). The \linearautomorphism \(f\) is
    \emph{orientation-reversing} if
    \[
        \det(f)<0.
    \]
\end{snippetdefinition}

\begin{snippetproposition}{automorphism-orientation-sign-dichotomy}{Automorphisms either preserve or reverse orientation}
    Let \(V\) be a finite-dimensional real \vectorspace, and let
    \(f\in\operatorname{Aut}(V)\). Then \(f\) is either
    \orientationpreserving or \orientationreversing.
\end{snippetproposition}

\begin{snippetproof}{automorphism-orientation-sign-dichotomy-proof}{automorphism-orientation-sign-dichotomy}{Automorphisms either preserve or reverse orientation}
    Since \(f\) is invertible, the associated \matrix of \(f\) in any
    \basis is invertible. Hence
    \[
        \det(f)\neq0.
    \]
    Over \(\realnumbers\), every nonzero scalar is either positive or
    negative, so either \(\det(f)>0\) or \(\det(f)<0\).
\end{snippetproof}

\section{Orientations of vector spaces}

\begin{snippetdefinition}{same-orientation-bases-definition}{Bases with the same orientation}
    Let \(V\) be a finite-dimensional real \vectorspace, and let
    \(\mathcal{B}\) and \(\mathcal{E}\) be two \basis[bases] of \(V\).
    They have the \emph{same orientation} if
    \[
        \det\bigl(M(\identityfunc_V,\mathcal{B},\mathcal{E})\bigr)>0.
    \]
\end{snippetdefinition}

\begin{snippetproposition}{same-orientation-bases-equivalence}{Same orientation is an equivalence relation}
    Let \(V\) be a finite-dimensional real \vectorspace. The relation
    ``having the same orientation'' is an equivalence relation on the \set of
    all ordered \basis[bases] of \(V\).
\end{snippetproposition}

\begin{snippetproof}{same-orientation-bases-equivalence-proof}{same-orientation-bases-equivalence}{Same orientation is an equivalence relation}
    The relation is reflexive because
    \[
        M(\identityfunc_V,\mathcal{B},\mathcal{B})=\identmatrix{n},
    \]
    and \(\det(\identmatrix{n})=1>0\).

    It is symmetric because
    \[
        M(\identityfunc_V,\mathcal{E},\mathcal{B})
        =
        M(\identityfunc_V,\mathcal{B},\mathcal{E})^{-1}.
    \]
    If the \determinant of a \matrix is positive, then the \determinant of
    its inverse is also positive.

    It is transitive because
    \[
        M(\identityfunc_V,\mathcal{B},\mathcal{D})
        =
        M(\identityfunc_V,\mathcal{E},\mathcal{D})
        M(\identityfunc_V,\mathcal{B},\mathcal{E}).
    \]
    If \(\mathcal{B}\) has the same orientation as \(\mathcal{E}\), and
    \(\mathcal{E}\) has the same orientation as \(\mathcal{D}\), then both
    determinants on the right are positive. Their product is positive, so
    \(\mathcal{B}\) has the same orientation as \(\mathcal{D}\).
\end{snippetproof}

\begin{snippetproposition}{ordered-bases-have-two-orientation-classes}{There are two orientation classes}
    Let \(V\) be a finite-dimensional real \vectorspace with
    \[
        \lineardim V=n\geq1.
    \]
    The relation of having the same orientation partitions the ordered
    \basis[bases] of \(V\) into exactly two equivalence classes.
\end{snippetproposition}

\begin{snippetproof}{ordered-bases-have-two-orientation-classes-proof}{ordered-bases-have-two-orientation-classes}{There are two orientation classes}
    Let
    \[
        \mathcal{B}=\{v_1,\ldots,v_n\}
    \]
    be an ordered \basis of \(V\), and put
    \[
        \mathcal{C}=\{-v_1,v_2,\ldots,v_n\}.
    \]
    The change-of-basis \matrix from \(\mathcal{B}\) to \(\mathcal{C}\) has
    determinant \(-1\), so \(\mathcal{B}\) and \(\mathcal{C}\) do not have
    the same orientation.

    If \(\mathcal{E}\) is any ordered \basis of \(V\), then
    \[
        \det\bigl(M(\identityfunc_V,\mathcal{B},\mathcal{E})\bigr)
    \]
    is nonzero. If it is positive, then \(\mathcal{E}\) is in the class of
    \(\mathcal{B}\). If it is negative, then
    \[
        \det\bigl(M(\identityfunc_V,\mathcal{C},\mathcal{E})\bigr)>0,
    \]
    so \(\mathcal{E}\) is in the class of \(\mathcal{C}\). Hence there are
    exactly two classes.
\end{snippetproof}

\begin{snippetdefinition}{vector-space-orientation-definition}{Orientation of a real vector space}
    Let \(V\) be a finite-dimensional real \vectorspace. An
    \emph{orientation} of \(V\) is a choice of one of the two equivalence
    classes of ordered \basis[bases] determined by the same-orientation
    relation.
\end{snippetdefinition}

\begin{snippetdefinition}{canonical-orientation-definition}{Canonical orientation of \(\mathbb{R}^n\)}
    The \emph{canonical orientation} of \(\realnumbers^n\) is the
    \orientation represented by the canonical ordered \basis
    \[
        \{e_1,\ldots,e_n\}.
    \]
\end{snippetdefinition}

\begin{snippetexample}{r3-canonical-orientation-example}{Orientation in \(\mathbb{R}^3\)}
    In \(\realnumbers^3\), the canonical \orientation is represented by the
    ordered \basis
    \[
        \{e_1,e_2,e_3\}.
    \]
    Replacing exactly one vector of this ordered \basis by its negative gives
    an ordered \basis with the opposite orientation.
\end{snippetexample}

\end{document}
